\chapter{Álgebra Conmutativa}

\emph{Nota:} este tema se ha reorganizado añadiendo, en la Sección~\ref{sec:preliminares}, todas las definiciones algebraicas previas (grupos, anillos, morfismos, cocientes, retículos, álgebras y producto tensorial) que se necesitan para entender con soltura el resto del capítulo. A partir de la Sección~\ref{sec:anillos-conmutativos} se sigue el mismo desarrollo que en clase, corrigiendo los errores y completando las demostraciones y definiciones que quedaron pendientes.

\section{Preliminares algebraicos}\label{sec:preliminares}

\subsection{Grupos abelianos}

\begin{definicion}
Un \textbf{grupo abeliano} es un par $(A,+)$ donde $A$ es un conjunto y
\Func{+}{A\times A}{A}{(a,b)}{a+b}
es una operación tal que:
\begin{enumerate}
    \item (Asociativa) $(a+b)+c = a+(b+c)$ para todo $a,b,c\in A$.
    \item (Conmutativa) $a+b=b+a$ para todo $a,b \in A$.
    \item (Elemento neutro) Existe $0_A\in A$ tal que $a+0_A = a$ para todo $a\in A$.
    \item (Opuesto) Para todo $a\in A$ existe $-a\in A$ tal que $a+(-a)=0_A$.
\end{enumerate}
\end{definicion}

Un \textbf{subgrupo} $H\leq A$ es un subconjunto no vacío cerrado por la suma y por opuestos ($0_A\in H$, $h,h'\in H \Rightarrow h+h'\in H$, $h\in H\Rightarrow -h\in H$). Como $A$ es abeliano, todo subgrupo es normal, así que siempre podremos formar el grupo cociente $A/H$.

\subsection{Anillos conmutativos}

\begin{definicion}\label{def:anillo}
Un \textbf{anillo conmutativo (con unidad)} es una terna $(A,+,\cdot)$ donde:
\begin{enumerate}
    \item $(A,+)$ es un grupo abeliano, con neutro $0_A$.
    \item La multiplicación
    \Func{\cdot}{A\times A}{A}{(a,a')}{a\cdot a'}
    es asociativa, conmutativa y posee elemento neutro $1_A$ (con $1_A\neq 0_A$, salvo que digamos explícitamente que $A$ es el \textbf{anillo nulo}).
    \item Se verifica la propiedad distributiva: $a\cdot(b+c) = a\cdot b + a\cdot c$ para todo $a,b,c\in A$.
\end{enumerate}
\end{definicion}

Un \textbf{subanillo} $B\subseteq A$ es un subconjunto que es anillo con las operaciones restringidas y con $1_A\in B$.

\begin{definicion}
Un elemento $a\in A$ es una \textbf{unidad} (o elemento invertible) si existe $b\in A$ con $ab=1_A$. El conjunto de unidades se denota $U(A)$ y forma un grupo abeliano con la multiplicación (el \textbf{grupo de unidades}).
\end{definicion}

Un anillo $A\neq 0$ es un \textbf{dominio de integridad} (o simplemente \textbf{dominio}) si no tiene divisores de cero: $ab=0 \Rightarrow a=0$ o $b=0$. Un anillo $A$ es un \textbf{cuerpo} si $A\neq 0$ y $U(A)=A\setminus\{0\}$, es decir, todo elemento no nulo es invertible. Todo cuerpo es un dominio.

\subsection{Morfismos de anillos}

\begin{definicion}
Sean $A,B$ anillos conmutativos. Un \textbf{morfismo de anillos} es una aplicación
\Func{\varphi}{A}{B}{a}{\varphi(a)}
tal que:
\begin{enumerate}
    \item $\varphi(a+a') = \varphi(a)+\varphi(a')$,
    \item $\varphi(a\cdot a') = \varphi(a)\cdot \varphi(a')$,
    \item $\varphi(1_A) = 1_B$.
\end{enumerate}
Se definen el \textbf{núcleo} y la \textbf{imagen} de $\varphi$ como
\begin{gather*}
    \ker(\varphi) = \{a\in A : \varphi(a)=0_B\}, \qquad
    \Im(\varphi) = \{b\in B : \exists a\in A,\ \varphi(a)=b\}.
\end{gather*}
\end{definicion}

Es inmediato comprobar que $\Im(\varphi)$ es un subanillo de $B$, y que $\ker(\varphi)$ es un \textbf{ideal} de $A$ (concepto que precisamos a continuación). Además $\varphi$ es inyectivo si y solo si $\ker(\varphi)=0$.

\begin{prop}
Para todo anillo conmutativo $A$ existe un único morfismo de anillos $\mathbb{Z}\to A$, dado por
\begin{gather*}
    n \mapsto \underbrace{1_A + \cdots + 1_A}_{n \text{ veces}} \quad (n\geq 0), \qquad
    n \mapsto -(-n)\cdot 1_A \quad (n<0).
\end{gather*}
Este morfismo es \textbf{inyectivo si y solo si $A$ tiene característica $0$} (es decir, $n\cdot 1_A\neq 0$ para todo $n\geq 1$); en general su núcleo es $n_0\mathbb{Z}$ donde $n_0=\mathrm{car}(A)\geq 0$ es la característica de $A$.
\end{prop}

\emph{(Corrección: en el enunciado original se afirmaba que este morfismo $\mathbb{Z}\to A$ es siempre inyectivo. Esto es falso en general: si $A=\mathbb{Z}_n$, el morfismo $\mathbb{Z}\to \mathbb{Z}_n$ tiene núcleo $n\mathbb{Z}\neq 0$. Lo que sí es cierto siempre es la existencia y unicidad del morfismo, ya que $\mathbb{Z}$ es el objeto inicial en la categoría de anillos conmutativos.)}

\subsection{Ideales}

\begin{definicion}
Sea $A$ un anillo conmutativo. Un subconjunto $I\subseteq A$ es un \textbf{ideal} de $A$, y escribimos $I\leq A$, si:
\begin{enumerate}
    \item $(I,+)\leq (A,+)$ es un subgrupo,
    \item $a\cdot x \in I$ para todo $a\in A$, $x\in I$ (propiedad de absorción).
\end{enumerate}
Un ideal $I$ es \textbf{propio} si $I\neq A$; equivalentemente, $1_A\notin I$.
\end{definicion}

Nótese que la condición de absorción se escribe de forma compacta como $A\cdot I \subseteq I$.

\subsection{Cociente de un conjunto por una relación de equivalencia}

Antes de cocientar anillos conviene recordar la construcción general. Sea $X$ un conjunto y $R$ una relación de equivalencia en $X$ (reflexiva, simétrica y transitiva). Para $x\in X$ definimos su \textbf{clase de equivalencia}
\begin{gather*}
    \bar{x} = \{x'\in X : xRx'\},
\end{gather*}
y el \textbf{conjunto cociente}
\begin{gather*}
    X/R = \{\bar{x} : x\in X\}.
\end{gather*}
Se tiene que $\bar x = \bar y$ si y solo si $xRy$, y la \textbf{proyección canónica}
\begin{gather*}
    \pi_R:X \to X/R, \qquad x\mapsto \bar x
\end{gather*}
es sobreyectiva.

\begin{prop}[Propiedad universal del cociente]\label{prop:universal-cociente}
Sea $f:X\to Y$ una aplicación tal que $xRx' \Rightarrow f(x)=f(x')$ (decimos que $f$ es \textbf{compatible con $R$}, o que \textbf{pasa al cociente}). Entonces existe una única aplicación $\bar f: X/R \to Y$ que hace conmutar el diagrama
\begin{gather*}
    \xymatrix{
        X \ar[d]_{\pi_R} \ar[r]^f & Y\\
        X/R \ar@{-->}[ur]_{\bar{f}}
    }
\end{gather*}
es decir, tal que $\bar f \circ \pi_R = f$. Concretamente $\bar f(\bar x) := f(x)$, y esto está bien definido precisamente por la hipótesis de compatibilidad.
\end{prop}

Esta propiedad universal es la que usaremos constantemente para comprobar que una operación "definida sobre representantes" no depende del representante elegido.

\subsection{Anillo cociente}

Sea $A$ un anillo conmutativo e $I\leq A$ un ideal. Como $(I,+)\leq (A,+)$, podemos considerar el grupo cociente $(A/I,+)$, cuyos elementos son las clases $\bar a = a+I$, con suma $\bar a + \bar a' = \overline{a+a'}$ (que está bien definida por ser $I$ subgrupo). Escribimos la proyección canónica como
\begin{gather*}
    \pi: A \to A/I \ \text{en } Ab, \qquad a\mapsto a+I=\bar a.
\end{gather*}

Queremos dotar a $A/I$ de una multiplicación, definida por $\bar a \cdot \bar{a'} := \overline{a\cdot a'}$. Veamos que está bien definida, es decir, que no depende de los representantes: si $\bar a=\bar b$ y $\bar{a'}=\bar{b'}$, entonces $a-b\in I$ y $a'-b'\in I$, y
\begin{gather*}
    aa'-bb' = a a' - b a' + ba' - bb' = (a-b)a' + b(a'-b') \in I
\end{gather*}
usando la absorción del ideal en ambos sumandos. Luego $\overline{aa'}=\overline{bb'}$, como queríamos.

Con esta multiplicación, $(A/I,+,\cdot)$ es un anillo conmutativo, y $\pi$ es un morfismo de anillos con $\ker(\pi)=I$.

\begin{prop}[Propiedad universal del anillo cociente]
Si $\varphi:A\to B$ es un morfismo de anillos con $I\subseteq \ker(\varphi)$, existe un único morfismo de anillos $\bar\varphi: A/I \to B$ tal que $\bar\varphi \circ \pi = \varphi$.
\end{prop}
\begin{proof}
Es la Proposición~\ref{prop:universal-cociente} aplicada a la relación de equivalencia $aRa' \sii a-a'\in I$, comprobando que $\bar\varphi$ así obtenida es morfismo de anillos (lo que se sigue de que $\varphi$ lo es).
\end{proof}

\subsection{Teoremas de isomorfía}

\begin{teo}[Primer teorema de isomorfía]
Si $\varphi:A\to B$ es un morfismo de anillos, entonces
\begin{gather*}
    A/\ker(\varphi) \cong \Im(\varphi).
\end{gather*}
\end{teo}

\begin{teo}[Teorema de correspondencia]\label{teo:correspondencia}
Sea $I\leq A$ y $\pi:A\to A/I$ la proyección canónica. La aplicación
\begin{gather*}
    \{J\leq A : I\subseteq J\} \longrightarrow \mathrm{Id}(A/I), \qquad J \mapsto J/I:=\pi(J)
\end{gather*}
es una biyección que respeta inclusiones (y por tanto un isomorfismo de retículos). Su inversa es $K\mapsto \pi^{-1}(K)$.
\end{teo}

\emph{Idea de la demostración.} Dado $K\leq A/I$ ideal, $J:=\pi^{-1}(K)=\{a\in A : a+I\in K\}$ es un ideal de $A$ que contiene a $I$ (pues $\bar 0 \in K$). Recíprocamente, si $I\subseteq J\leq A$, entonces $J/I=\pi(J)$ es un subgrupo de $A/I$, y es ideal porque para $\bar a\in A/I$, $\bar j\in J/I$:
\begin{gather*}
    \bar a \cdot \bar j = \overline{aj} \in J/I,
\end{gather*}
ya que $J$ es ideal de $A$ ($aj\in J$ para $a\in A,\ j\in J$). Se comprueba sin dificultad que ambas construcciones son inversas entre sí.

\begin{teo}[Tercer teorema de isomorfía]\label{teo:tercer-isomorfia}
Si $I\leq J \leq A$ son ideales, entonces $J/I \leq A/I$ es un ideal y
\begin{gather*}
    \frac{A/I}{J/I} \cong A/J.
\end{gather*}
\end{teo}
\begin{proof}
Consideramos el morfismo compuesto $A \xrightarrow{\pi} A/I \to (A/I)/(J/I)$, que tiene núcleo exactamente $J$ (un elemento $a$ va a $\bar 0$ si y solo si $\bar a \in J/I$, si y solo si $a\in \pi^{-1}(J/I)=J$ por el Teorema~\ref{teo:correspondencia}). Aplicando el primer teorema de isomorfía a este morfismo se obtiene el resultado.
\end{proof}

\subsection{Conjuntos parcialmente ordenados y retículos}

\begin{definicion}
Un \textbf{conjunto parcialmente ordenado} (\textbf{poset}) es un par $(P,\leq)$ donde $\leq$ es una relación reflexiva, antisimétrica y transitiva. Una \textbf{cadena} en $P$ es un subconjunto totalmente ordenado (cualesquiera dos elementos son comparables).
\end{definicion}

\begin{definicion}
Un \textbf{retículo} es un poset $(P,\leq)$ en el que cualesquiera dos elementos $x,y$ poseen:
\begin{itemize}
    \item un \textbf{ínfimo} (o encuentro) $x\wedge y$: el mayor elemento que es $\leq x$ y $\leq y$;
    \item un \textbf{supremo} (o unión) $x\vee y$: el menor elemento que es $\geq x$ y $\geq y$.
\end{itemize}
\end{definicion}

Más adelante veremos que el conjunto $\mathrm{Id}(A)$ de ideales de un anillo, ordenado por inclusión, es un retículo.

\subsection{$A$-álgebras}

\begin{definicion}
Sean $A,B$ anillos conmutativos. Decimos que $B$ es una \textbf{$A$-álgebra} si existe un morfismo de anillos $\varphi:A\to B$ (que llamamos \textbf{estructura de $A$-álgebra}). Un \textbf{morfismo de $A$-álgebras} $B\to C$ es un morfismo de anillos que es compatible con las estructuras, es decir, que hace conmutar el triángulo con $A$.
\end{definicion}

Como todo anillo conmutativo $A$ recibe un único morfismo $\mathbb{Z}\to A$, \textbf{todo anillo conmutativo es de forma única una $\mathbb{Z}$-álgebra}. Este es el motivo por el que muchas construcciones "sobre un cuerpo $\mathbb{K}$" se pueden repetir "sobre $\mathbb{Z}$" para un anillo cualquiera.

\subsection{Producto tensorial}

Necesitamos ahora el producto tensorial, que usaremos primero para $\mathbb{K}$-espacios vectoriales y después, de forma completamente análoga, para $A$-módulos ($\mathbb{Z}$-módulos en particular).

\begin{definicion}[Aplicación bilineal]
Sean $V,W,U$ espacios vectoriales sobre un cuerpo $\mathbb{K}$ (o módulos sobre un anillo $A$). Una aplicación $b:V\times W\to U$ es \textbf{bilineal} si es lineal en cada variable por separado.
\end{definicion}

\begin{definicion}[Producto tensorial]
El \textbf{producto tensorial} $V\otimes_{\mathbb{K}} W$ es un $\mathbb{K}$-espacio vectorial junto con una aplicación bilineal
\begin{gather*}
    \otimes: V\times W \to V\otimes_{\mathbb{K}} W, \qquad (v,w)\mapsto v\otimes w
\end{gather*}
que satisface la siguiente \textbf{propiedad universal}: para toda aplicación bilineal $b:V\times W\to U$ existe una única aplicación lineal $\tilde b: V\otimes_{\mathbb{K}} W \to U$ tal que $\tilde b(v\otimes w)=b(v,w)$ para todo $v,w$.
\end{definicion}

Esta propiedad determina a $V\otimes_{\mathbb{K}} W$ de forma única salvo isomorfismo. Una construcción explícita: se toma el espacio vectorial libre con base $V\times W$ y se cocienta por el subespacio generado por los elementos que expresan la bilinealidad, es decir, por
\begin{gather*}
    (v+v',w)-(v,w)-(v',w),\quad (v,w+w')-(v,w)-(v,w'),\quad (\lambda v, w)-\lambda(v,w),\quad (v,\lambda w)-\lambda(v,w).
\end{gather*}
La misma construcción, sustituyendo $\mathbb{K}$ por un anillo $A$ y "espacio vectorial" por "$A$-módulo", define $M\otimes_A N$ para $A$-módulos $M,N$ cualesquiera.

\medskip
\textbf{Propiedades básicas} (para $V,W$ de dimensión finita sobre $\mathbb{K}$):
\begin{itemize}
    \item $\dim_{\mathbb{K}}(V\otimes_{\mathbb{K}} W) = \dim(V)\cdot \dim(W)$.
    \item Si $\{e_i\}$ es base de $V$ y $\{f_j\}$ base de $W$, entonces $\{e_i\otimes f_j\}$ es base de $V\otimes_{\mathbb{K}} W$.
    \item Extensión de escalares: si $A$ es una $\mathbb{K}$-álgebra, entonces $A\otimes_{\mathbb{K}} \mathbb{K}[X] \cong A[X]$, y en general, para cualquier anillo conmutativo $A$ (visto como $\mathbb{Z}$-álgebra),
    \begin{gather*}
        A[X] \cong A \underset{\mathbb{Z}}{\otimes} \mathbb{Z}[X].
    \end{gather*}
\end{itemize}

\emph{(Corrección: en el diagrama original aparecía $A\otimes_{\mathbb{K}}\mathbb{Z}[X]$, mezclando un producto tensorial sobre $\mathbb{K}$ con el anillo $\mathbb{Z}[X]$, que en general no es un $\mathbb{K}$-espacio vectorial. La identidad válida en general es sobre $\mathbb{Z}$; solo si $A$ es ya una $\mathbb{K}$-álgebra tiene sentido escribir $A\otimes_{\mathbb{K}} \mathbb{K}[X]\cong A[X]$.)}

\subsubsection{Dual de un espacio vectorial y $\mathrm{Hom}$}

Si $V$ es un $\mathbb{K}$-espacio vectorial de dimensión finita, su \textbf{dual} es
\begin{gather*}
    V^* = \mathrm{Hom}_{\mathbb{K}}(V,\mathbb{K}),
\end{gather*}
el espacio de aplicaciones lineales de $V$ en $\mathbb{K}$, con $\dim(V^*)=\dim(V)$.

Si $\dim(V)=n$ y $\dim(W)=m$ son finitas, fijando bases se tiene
\begin{gather*}
    \dim(\mathrm{Hom}_{\mathbb{K}}(V,W)) = \dim(M_{m\times n}(\mathbb{K})) = n\cdot m,
\end{gather*}
(cada aplicación lineal $V\to W$ se representa por una matriz $m\times n$), y hay un isomorfismo canónico
\begin{gather*}
    V^* \otimes_{\mathbb{K}} W \;\cong\; \mathrm{Hom}_{\mathbb{K}}(V,W), \qquad \varphi \otimes w \mapsto \big(v\mapsto \varphi(v)\, w\big),
\end{gather*}
consistente con $\dim(V^*\otimes W) = \dim(V^*)\cdot \dim(W) = n\cdot m$.

\subsubsection{Sumas y productos directos; soporte finito}

Dada una familia $\{W_n\}_{n\in \mathbb{N}}$ de espacios vectoriales, el \textbf{producto directo} $\prod_{n\in \mathbb{N}} W_n$ es el conjunto de todas las sucesiones $(x_n)_{n\in \mathbb{N}}$ con $x_n\in W_n$. La \textbf{suma directa}
\begin{gather*}
    \bigoplus_{n\in \mathbb{N}} W_n \;\subseteq\; \prod_{n\in \mathbb{N}} W_n
\end{gather*}
es el subconjunto de las sucesiones de \textbf{soporte finito}, es decir, aquellas $(x_n)_n$ con $x_n=0$ salvo para un número finito de índices $n$.

\subsubsection{Álgebra tensorial, simétrica y exterior}

Sea $W$ un $\mathbb{K}$-espacio vectorial. Para $n\geq 0$ escribimos $W^{\otimes n} = W\otimes_{\mathbb{K}} \cdots \otimes_{\mathbb{K}} W$ ($n$ veces), con el convenio $W^{\otimes 0}=\mathbb{K}$. El \textbf{álgebra tensorial} de $W$ es la suma directa (de soporte finito, como arriba)
\begin{gather*}
    T(W) = \bigoplus_{n\geq 0} W^{\otimes n},
\end{gather*}
con producto dado por la concatenación de tensores:
\begin{gather*}
    (v_1\otimes\cdots\otimes v_k)\cdot (w_1\otimes \cdots \otimes w_l) := v_1\otimes\cdots\otimes v_k \otimes w_1 \otimes\cdots\otimes w_l.
\end{gather*}
Un elemento genérico de $T(W)$ es entonces una suma finita de tensores puros de distintos órdenes, por ejemplo
\begin{gather*}
    \lambda_0 \cdot 1 + w_1 + w_2\otimes w_3 + \cdots + w_{k_1}\otimes \cdots \otimes w_{k_r} ,
\end{gather*}
donde solo un número finito de sumandos son no nulos.

A partir de $T(W)$ se construyen dos cocientes muy importantes, tomando el ideal (bilátero) generado por ciertas relaciones dentro de $T(W)$:

\begin{itemize}
    \item El \textbf{álgebra simétrica}
    \begin{gather*}
        S(W) = T(W)/\langle v\otimes w - w\otimes v \ :\ v,w\in W\rangle,
    \end{gather*}
    en la que el producto se vuelve conmutativo. De hecho $S(W)\cong \mathbb{K}[x_1,\ldots,x_n]$ si $\dim(W)=n$.
    \item El \textbf{álgebra exterior}
    \begin{gather*}
        \Lambda(W) = T(W)/\langle v\otimes v \ : \ v\in W\rangle,
    \end{gather*}
    cuyo producto se denota $v\wedge w$ y es \textbf{anticonmutativo}: $v\wedge w = -w\wedge v$ (lo cual se deduce de imponer $v\otimes v=0$: en efecto $0=(v+w)\otimes(v+w) = v\otimes v + v\otimes w + w\otimes v + w\otimes w = v\otimes w+w\otimes v$, de donde $v\wedge w=-w\wedge v$; si $\mathrm{car}(\mathbb{K})\neq 2$ esto equivale a imponer directamente $v\otimes w+w\otimes v=0$).
\end{itemize}

\emph{(Corrección: la fórmula original, $I=\langle v\otimes w+w\otimes v\rangle$ seguida de "$\wedge W = T_{\mathbb{K}}$", tenía dos problemas: faltaba el cociente ("$/I$") y la notación $T_{\mathbb{K}}$ no correspondía a nada definido. La definición correcta y válida en toda característica usa el generador $v\otimes v$, como arriba.)}

\bigskip
Con todos estos preliminares ya podemos leer el resto del tema con rigor.

\section{Anillos conmutativos: primeros ejemplos}\label{sec:anillos-conmutativos}

Recordamos (Definición~\ref{def:anillo}) que un anillo conmutativo es un grupo abeliano $(A,+)$ con una multiplicación asociativa, conmutativa, con neutro $1_A$, distributiva respecto de la suma.

\begin{ejemplo}\
    \begin{itemize}
        \item El primer ejemplo que vamos a tener en cuenta es la cadena de inclusiones de cuerpos
        \begin{gather*}
            \mathbb{Z}\subset \mathbb{Q} \subset \mathbb{R} \subset \mathbb{C} \subset \mathbb{H} \subset \mathbb{O},
        \end{gather*}
        donde $\mathbb{H}$ son los cuaterniones y $\mathbb{O}$ los octoniones, obtenidos a partir de $\mathbb{C}$ mediante la \textbf{construcción de Cayley-Dickson}. En cada paso se va perdiendo alguna propiedad: $\mathbb{H}$ ya no es conmutativo (por lo que no es, en el sentido de este tema, un anillo conmutativo, aunque sí un anillo con división) y $\mathbb{O}$ ni siquiera es asociativo. Estos son ejemplos clásicos y a partir de ellos podremos construir más.

        \emph{(Corrección: la nota original escribía $\mathbb{R}=\bar{\mathbb{Q}}$, es decir, identificaba $\mathbb{R}$ con la clausura algebraica de $\mathbb{Q}$. Esto es falso: la clausura algebraica de $\mathbb{Q}$ dentro de $\mathbb{C}$, que se suele denotar $\bar{\mathbb{Q}}$, es el cuerpo (numerable) de los números algebraicos, estrictamente más pequeño que $\mathbb{R}$, que no es algebraicamente cerrado ni numerable.)}

        \item Consideramos un cuerpo $\mathbb{K}$ y un conjunto $X\neq \emptyset$. El conjunto $\mathrm{Maps}(X,\mathbb{K})$ de todas las aplicaciones de $X$ en $\mathbb{K}$, con las operaciones puntuales, es un $\mathbb{K}$-álgebra, y es naturalmente isomorfo al producto
        \begin{gather*}
            \mathbb{K}^X = \prod_{x\in X} \mathbb{K}
        \end{gather*}
        (a cada función $f$ le asociamos la tupla $(f(x))_{x\in X}$).

        Tomemos ahora $X=\mathbb{N}$. Dentro de $\mathrm{Maps}(\mathbb{N},\mathbb{K})=\mathbb{K}^{\mathbb{N}}$ (todas las sucesiones) tenemos el subconjunto $\mathbb{K}^{(\mathbb{N})}$ de \textbf{soporte finito} (sucesiones casi nulas), que es un ideal de $\mathbb{K}^{\mathbb{N}}$ pero \emph{no} contiene a la unidad, así que como anillo se estudia aparte. Sobre las sucesiones (finitas o de soporte finito) de escalares podemos considerar dos productos distintos:
        \begin{itemize}
            \item el \textbf{producto de Cauchy} (el producto usual de series/polinomios de potencias),
            \begin{gather*}
                (a_n)_n * (b_n)_n = \Big(\sum_{k=0}^{n} a_k\, b_{n-k}\Big)_n,
            \end{gather*}
            \item el \textbf{producto de Hurwitz}, que incorpora coeficientes binomiales,
            \begin{gather*}
                (a_n)_n \cdot (b_n)_n = \Big(\sum_{k=0}^{n} \binom{n}{k}\, a_k\, b_{n-k}\Big)_n.
            \end{gather*}
        \end{itemize}
        (\emph{Corrección}: en el enunciado original aparecía $b_{n\cdot k}$, un error tipográfico; el índice correcto es $b_{n-k}$.) El producto de Hurwitz corresponde a la estructura de anillo en las series de potencias exponenciales (útil, por ejemplo, al trabajar con funciones generatrices exponenciales).

        \item Si $A$ es un anillo conmutativo cualquiera, existe un único morfismo $\mathbb{Z}\to A$ (visto arriba), lo que hace de $A$ una $\mathbb{Z}$-álgebra de forma canónica. Esto nos permite repetir muchas construcciones "sobre $\mathbb{K}$" trabajando "sobre $\mathbb{Z}$" para cualquier $A$.

        Recordamos: $B$ es una $A$-álgebra si existe un morfismo de anillos $A\to B$.

        Se tiene entonces la identidad (extensión de escalares, vista en los preliminares)
        \begin{gather*}
            A[X] \;\cong\; A\underset{\mathbb{Z}}{\otimes} \mathbb{Z}[X].
        \end{gather*}

        \item Si $V$ es un $\mathbb{K}$-espacio vectorial con $\dim_{\mathbb{K}}(V)<\infty$, su dual es $V^* = \mathrm{Hom}(V,\mathbb{K})$ (visto arriba).
    \end{itemize}
\end{ejemplo}

Como se explicó en los preliminares, si $W$ es un $\mathbb{K}$-espacio vectorial podemos formar $W^{\otimes n}$, el producto $\prod_{n\geq 0} W^{\otimes n}$ y, dentro de él, la suma directa de soporte finito $T(W)=\bigoplus_{n\geq 0} W^{\otimes n}$ (el álgebra tensorial), con el producto dado por concatenación. Recordamos también que si $\dim(V),\dim(W)<\infty$,
\begin{gather*}
    \dim(\mathrm{Hom}(V,W)) = \dim(M_{n\times m}(\mathbb{K})) = n\cdot m = \dim(V^*\otimes W),
\end{gather*}
lo que es consistente con el isomorfismo $\mathrm{Hom}(V,W)\cong V^*\otimes W$.

Finalmente, cocientando $T(W)$ por el ideal $\langle v\otimes v\rangle_{v\in W}$ obtenemos el álgebra exterior $\Lambda(W)$, con producto anticonmutativo $v\wedge w=-w\wedge v$; cocientando en cambio por $\langle v\otimes w - w\otimes v\rangle$ obtenemos el álgebra simétrica $S(W)$, conmutativa.

\section{Ideales, morfismos y anillo cociente}

Recogemos aquí, aplicadas ya al caso de anillos, las nociones generales de los preliminares. Un morfismo de anillos $A\xrightarrow{\varphi} B$ tiene
\begin{enumerate}
    \item $\ker(\varphi) = \{a \in A : \varphi(a) = 0\} \leq A$ (ideal),
    \item $\Im(\varphi) = \{b\in B : \exists a \in A,\ \varphi(a) = b\} \leq B$ (subanillo).
\end{enumerate}
Dado $A$ un anillo e $I\leq A$ un ideal, tenemos el subgrupo $(I,+) \leq (A,+)$, la absorción $A \cdot I = I \cdot A \subseteq I$, y la proyección canónica en la categoría de grupos abelianos
\begin{gather*}
    A \xrightarrow{\pi} A/I \text{ en } Ab, \qquad a \mapsto a + I = \bar{a},
\end{gather*}
que, como vimos, se puede dotar de estructura de anillo mediante $\bar{a} \cdot \bar{a'} := \overline{aa'}$, una definición que está bien planteada (no depende de representantes) exactamente por el cálculo hecho en la Sección~\ref{sec:preliminares}.

La propiedad universal del cociente (Prop.~\ref{prop:universal-cociente}, aplicada a la relación $aRa'\sii a-a'\in I$) es la herramienta que sistematiza todo esto: si $f:A\to Y$ es compatible con la relación de congruencia módulo $I$, factoriza de forma única a través de $\pi$.

Con esta maquinaria podemos hacer muchas cosas. Por ejemplo, dada una torre de ideales $I \leq J \leq A$, el Teorema~\ref{teo:correspondencia} nos dice que los ideales de $A/I$ son exactamente los $J/I$ con $I\leq J\leq A$, y el Teorema~\ref{teo:tercer-isomorfia} nos da
\begin{gather*}
    \frac{A/I}{J/I} \;\cong\; A/J.
\end{gather*}

Veamos con más detalle por qué $J/I$ es efectivamente un ideal de $A/I$: es un subgrupo (imagen de un subgrupo por un morfismo de grupos), y para $\alpha=\bar a \in A/I$ y $x=\bar j\in J/I$ (con $j\in J$),
\begin{gather*}
    x\cdot \alpha = \bar{j}\cdot \bar{a} = \overline{ja} \in J/I,
\end{gather*}
ya que $ja \in J$ por ser $J$ ideal de $A$.

\begin{ejemplo}
Sabemos que $\mathbb{Z}_n=\mathbb{Z}/n\mathbb{Z}$ son anillos finitos para $n>1$. Comparemos, para $n=12$, el conjunto $D(12)=\{1,2,3,4,6,12\}$ de divisores positivos de $12$, ordenado por divisibilidad, con el retículo $\mathrm{Id}(\mathbb{Z}_{12})$ de ideales de $\mathbb{Z}_{12}$.

Por el Teorema~\ref{teo:correspondencia}, los ideales de $\mathbb{Z}_{12}=\mathbb{Z}/12\mathbb{Z}$ son exactamente los $d\mathbb{Z}/12\mathbb{Z} = \langle \bar d\rangle$ para $d$ divisor de $12$, es decir hay exactamente $6$ ideales, uno por cada $d\in D(12)$:
\begin{gather*}
    \langle \bar 1\rangle = \mathbb{Z}_{12}, \quad \langle \bar 2\rangle,\quad \langle \bar 3\rangle,\quad \langle \bar 4\rangle,\quad \langle \bar 6\rangle,\quad \langle\bar{12}\rangle=\langle \bar 0\rangle=0.
\end{gather*}
La aplicación $d\mapsto \langle \bar d\rangle$ es una biyección entre $D(12)$ y $\mathrm{Id}(\mathbb{Z}_{12})$ que \textbf{invierte} el orden: $d\,|\,d' \iff \langle \bar{d'}\rangle \subseteq \langle \bar d\rangle$ (cuanto mayor es el divisor, más pequeño es el ideal que genera). Así, el retículo de ideales de $\mathbb{Z}_{12}$ es (isomorfo a) el retículo dual de $(D(12), \,|\,)$.
\end{ejemplo}

\section{El retículo de ideales}

\begin{prop}
Para cualquier anillo conmutativo $A$, el conjunto $\mathrm{Id}(A)$ de ideales de $A$, ordenado por inclusión, es un retículo, con
\begin{gather*}
    I \wedge J = I \cap J, \qquad\qquad I \vee J = I + J.
\end{gather*}
\end{prop}

\emph{(Corrección: en el original se escribía $I\vee J = A+J$, que no tiene sentido —$A+J=A$ siempre—; la operación correcta es $I+J = \{a+b: a\in I, b\in J\}$.)}

Para $S \subseteq A$ un subconjunto cualquiera (no necesariamente un ideal), definimos el \textbf{ideal generado} por $S$ como el menor ideal que contiene a $S$:
\begin{gather*}
    \langle S \rangle = \bigcap_{S \subseteq I \leq A} I = \Big\{\sum_{\text{finita}} a_i s_i \ :\ a_i\in A,\ s_i\in S\Big\}.
\end{gather*}
(La segunda igualdad se comprueba viendo que el conjunto de sumas finitas de la derecha es ya un ideal que contiene a $S$, y que está contenido en cualquier ideal que contenga a $S$.)

Más generalmente, dada una familia $\{I_\lambda\}_{\lambda \in \Lambda}$ de ideales,
\begin{gather*}
    \sum_{\lambda \in \Lambda} I_\lambda := \Big\langle \bigcup_{\lambda \in \Lambda} I_\lambda \Big\rangle = \Big\{\sum_{\text{finita}} x_\lambda \ : \ x_\lambda \in I_\lambda\Big\},
\end{gather*}
es el ideal generado por la unión (nótese que la unión de ideales no es, en general, un ideal, pero sí lo es la unión de una \emph{cadena} de ideales, hecho que usaremos en la siguiente sección).

\section{Ideales primos y maximales. Teorema de Krull}

\begin{definicion}
Sea $A$ un anillo conmutativo. Un ideal $M\leq A$ es \textbf{maximal} si $M\neq A$ y no existe ningún ideal $J$ con $M\subsetneq J \subsetneq A$. Escribimos $\mathrm{Max}(A)$ para el conjunto de ideales maximales.
\end{definicion}

\begin{teo}[Lema de Zorn]\label{lema:zorn}
Sea $(P,\leq)$ un conjunto parcialmente ordenado no vacío en el que toda cadena posee una cota superior en $P$. Entonces $P$ posee (al menos) un elemento maximal.
\end{teo}

\emph{(Precisión: el Lema de Zorn exige que toda cadena tenga una \textbf{cota superior} en $P$, no necesariamente un \textbf{supremo}; es una diferencia sutil pero importante —una cota superior no tiene por qué ser la menor posible—. El Lema de Zorn es, de hecho, equivalente al Axioma de Elección, y no se demuestra a partir de los axiomas de Zermelo-Fraenkel sin él; lo tomamos como axioma.)}

\begin{teo}[Teorema de Krull]\label{teo:krull}
Sea $A$ un anillo conmutativo no nulo, $I\leq A$ un ideal propio ($I\neq A$; se permite $I=0$). Entonces existe $M\in \mathrm{Max}(A)$ tal que $I \subseteq M$. En particular, $\mathrm{Max}(A)\neq\emptyset$.
\end{teo}
\begin{proof}
Consideramos el conjunto
\begin{gather*}
    P_I = \{J \leq A : J \text{ ideal propio},\ I \subseteq J\},
\end{gather*}
ordenado por inclusión. Es no vacío porque $I\in P_I$. Sea $\{J_\lambda\}_{\lambda\in \Lambda}$ una cadena en $P_I$ y consideremos $J=\bigcup_{\lambda} J_\lambda$. Como es unión de una cadena de ideales (dos elementos cualesquiera de $J$ están, cada uno, en algún $J_\lambda$, y como la cadena está totalmente ordenada, ambos están en el mayor de los dos), $J$ es un ideal. Además $J$ es propio: si $1_A\in J$, entonces $1_A \in J_\lambda$ para algún $\lambda$, contradiciendo que $J_\lambda$ es propio. Así $J\in P_I$ y es cota superior de la cadena.

Por el Lema de Zorn (Teorema~\ref{lema:zorn}), $P_I$ tiene un elemento maximal $M$. Este $M$ es maximal también en $\mathrm{Id}(A)$ (no solo en $P_I$): si $M \subsetneq J \subsetneq A$ para algún ideal $J$, entonces $J\in P_I$ (pues $I\subseteq M\subseteq J$) y $J\supsetneq M$, contradiciendo la maximalidad de $M$ en $P_I$. Luego $M\in \mathrm{Max}(A)$ e $I\subseteq M$.
\end{proof}

\begin{definicion}
Un ideal propio $\mathfrak p \leq A$ es \textbf{primo} si
\begin{gather*}
    ab\in \mathfrak p \ \Rightarrow \ a\in \mathfrak p \text{ o } b\in \mathfrak p.
\end{gather*}
Escribimos $\mathrm{Spec}(A) = \{\mathfrak p \in \mathrm{Id}(A) : \mathfrak p \text{ primo}\}$, el \textbf{espectro} de $A$.
\end{definicion}

\begin{prop}\label{prop:primo-equivalencias}
Para un ideal propio $\mathfrak p \leq A$, son equivalentes:
\begin{enumerate}
    \item $\mathfrak p$ es primo.
    \item $A/\mathfrak p$ es un dominio de integridad.
    \item $A\setminus \mathfrak p$ es un subconjunto multiplicativamente cerrado (contiene a $1$ y es cerrado por productos).
    \item Para todo par de ideales $I,J\leq A$: $IJ\subseteq \mathfrak p \Rightarrow I\subseteq \mathfrak p$ o $J\subseteq \mathfrak p$.
\end{enumerate}
\end{prop}
\begin{proof}
(1)$\Leftrightarrow$(2)$\Leftrightarrow$(3) son reformulaciones directas de la definición de "$\mathfrak p$ primo" pasando al cociente: $ab\in\mathfrak p \Leftrightarrow \bar a\bar b=0$ en $A/\mathfrak p$, y "$a\notin \mathfrak p$ o $b\notin \mathfrak p$" es "$\bar a\neq 0$ o $\bar b\neq 0$"; la condición se traduce en que $A/\mathfrak p$ no tiene divisores de cero, que es (2), y (3) es la contrarrecíproca de (2) leída en $A$.

(1)$\Rightarrow$(4): supongamos $IJ\subseteq \mathfrak p$ e $I\not\subseteq \mathfrak p$. Tomamos $a\in I\setminus \mathfrak p$. Para todo $b\in J$, $ab\in IJ\subseteq \mathfrak p$; como $a\notin \mathfrak p$ y $\mathfrak p$ es primo, $b\in \mathfrak p$. Luego $J\subseteq \mathfrak p$.

(4)$\Rightarrow$(1): dados $a,b$ con $ab\in \mathfrak p$, tomamos $I=\langle a\rangle$, $J=\langle b\rangle$; entonces $IJ=\langle ab\rangle \subseteq \mathfrak p$, y por (4), $I\subseteq \mathfrak p$ o $J\subseteq \mathfrak p$, es decir $a\in \mathfrak p$ o $b\in\mathfrak p$.
\end{proof}

\emph{(Corrección importante: en los apuntes originales se añadía una tercera propiedad supuestamente equivalente, "$I\cap J\subseteq \mathfrak p \Rightarrow I\subseteq \mathfrak p$ o $J\subseteq \mathfrak p$". Esta propiedad \textbf{sí se deduce} de ser primo —porque $IJ\subseteq I\cap J$, así que $I\cap J\subseteq \mathfrak p$ implica $IJ\subseteq \mathfrak p$ y se aplica (4)— pero \textbf{no es equivalente}: hay ideales no primos que la verifican. Por ejemplo, en $\mathbb{Z}$, el ideal $\mathfrak p=(4)$ no es primo ($2\cdot 2\in (4)$ pero $2\notin (4)$); sin embargo, si $I=(m)$, $J=(n)$ son tales que $(m)\cap(n)=(\mathrm{mcm}(m,n))\subseteq (4)$, entonces $4\mid \mathrm{mcm}(m,n)$, lo cual fuerza a que $4\mid m$ o $4\mid n$ (si ambos tuvieran a lo sumo un factor $2$, también lo tendría su mcm), es decir $I\subseteq (4)$ o $J\subseteq (4)$. Así, $(4)$ verifica la propiedad de la intersección sin ser primo: la propiedad "de la intersección" es, en general, estrictamente más débil que ser primo.)}

\begin{definicion}
Un ideal propio $P\leq A$ es \textbf{irreducible} si no puede escribirse como $P=I\cap J$ con $I,J$ ideales que contienen \emph{estrictamente} a $P$. Escribimos $\mathrm{Irr}(A)$ para el conjunto de ideales irreducibles.
\end{definicion}

\begin{prop}
$\mathrm{Max}(A) \subseteq \mathrm{Spec}(A) \subseteq \mathrm{Irr}(A)$.
\end{prop}
\begin{proof}
$\mathrm{Max}(A)\subseteq \mathrm{Spec}(A)$: si $M$ es maximal, $A/M$ es un anillo no nulo en el que el único ideal propio es $0$ (por el Teorema~\ref{teo:correspondencia}, no hay ideales entre $M$ y $A$), luego todo elemento no nulo de $A/M$ es invertible (genera el ideal total, que es $A/M$), es decir $A/M$ es un cuerpo; en particular es un dominio, así que $M$ es primo por la Proposición~\ref{prop:primo-equivalencias}.

$\mathrm{Spec}(A)\subseteq \mathrm{Irr}(A)$: sea $\mathfrak p$ primo y supongamos $\mathfrak p = I\cap J$ con $I,J\supsetneq \mathfrak p$. Como $IJ\subseteq I\cap J=\mathfrak p$, por (4) de la Proposición~\ref{prop:primo-equivalencias}, $I\subseteq \mathfrak p$ o $J\subseteq \mathfrak p$; combinado con $I,J\supseteq \mathfrak p$ tendríamos $I=\mathfrak p$ o $J=\mathfrak p$, contradiciendo que ambos son estrictos. Luego $\mathfrak p$ es irreducible.
\end{proof}

\begin{lema}
Sea $x\in U(A)$. Es equivalente a decir que $x\notin M$ para todo $M\in \mathrm{Max}(A)$.
\end{lema}
\begin{proof}\
\begin{description}
    \item[$\Rightarrow$)] Trivial: si $x$ es unidad y $x\in M$ para algún ideal propio $M$, entonces $1=x^{-1}x\in M$, luego $M=A$, contradicción.
    \item[$\Leftarrow$)] Supongamos $x\notin M$ para todo $M\in \mathrm{Max}(A)$, y veamos que $x$ es unidad. Si $x=0$, como $A\neq 0$ hay algún maximal (Krull aplicado a $I=0$) y $0$ estaría en él, contradicción; así que $x\neq 0$. Supongamos, para llegar a una contradicción, que $x\notin U(A)$. Entonces el ideal $\langle x\rangle$ es propio (si $\langle x\rangle=A$, entonces $1\in \langle x\rangle$, es decir $1=ax$ para algún $a$, y $x$ sería unidad). Por el Teorema de Krull (Teorema~\ref{teo:krull}), existe $M\in \mathrm{Max}(A)$ con $\langle x\rangle \subseteq M$, luego $x\in M$, contradiciendo la hipótesis.
\end{description}
\end{proof}

\emph{(Corrección: en el enunciado original se escribía "$x\in I(A)$", una notación no definida; el resultado correcto y el que se usa después —de hecho, para \textbf{caracterizar los anillos locales}— es sobre el grupo de unidades $U(A)$.)}

\section{Anillos locales}

\begin{definicion}
Un anillo $A\neq 0$ se dice \textbf{local} si posee un único ideal maximal, es decir, $\mathrm{Max}(A)=\{M\}$ para cierto ideal $M$.
\end{definicion}

\begin{prop}
$A$ es local si y solo si $A\setminus U(A)$ es un ideal. En tal caso, $M=A\setminus U(A)$ es el único ideal maximal.
\end{prop}
\begin{proof}\
\begin{description}
    \item[$\Rightarrow$)] Supongamos $A$ local, con único maximal $M$. Veamos $A\setminus U(A)=M$.

    Si $x\in M$, entonces $x\notin U(A)$ (si $x$ fuera unidad, $M=\langle x\rangle_{\text{ideal generado}}\cdot A \ni 1$, imposible en un ideal propio). Así $M\subseteq A\setminus U(A)$.

    Recíprocamente, si $x\notin U(A)$, por el lema anterior existe algún maximal que contiene a $x$; como $A$ es local, ese maximal es $M$, luego $x\in M$. Así $A\setminus U(A)\subseteq M$.

    Por tanto $A\setminus U(A)=M$, que es un ideal.

    \item[$\Leftarrow$)] Supongamos que $M:=A\setminus U(A)$ es un ideal. Es propio, pues $1\in U(A)$, luego $1\notin M$. Veamos que es el único maximal (y por tanto que $A$ es local).

    Primero, $M$ es maximal: si $M\subsetneq J \leq A$, tomamos $x\in J\setminus M$; como $x\notin M=A\setminus U(A)$, $x$ es unidad, luego $J\ni x^{-1}x=1$, así $J=A$. Esto muestra que no hay ideales propios estrictamente mayores que $M$, es decir, $M$ es maximal.

    Ahora, si $M'$ es cualquier ideal maximal, en particular es propio, así que $M'\cap U(A)=\emptyset$ (ningún ideal propio contiene unidades), es decir $M'\subseteq A\setminus U(A)=M$. Como $M'$ es maximal y $M\neq A$, forzosamente $M'=M$.

    Luego $M$ es el único ideal maximal.
\end{description}
\end{proof}

\section{$\mathbb{R}$-álgebras, el espectro maximal y funciones continuas}

Consideramos ahora $\mathbb{K}=\mathbb{R}$ y $A$ una $\mathbb{R}$-álgebra (es decir, un anillo conmutativo con un morfismo $\mathbb{R}\to A$, en particular una estructura de $\mathbb{R}$-espacio vectorial compatible con el producto). Denotamos por
\begin{gather*}
    \mathrm{Alg}_{\mathbb{R}}(A,\mathbb{R})
\end{gather*}
el conjunto de morfismos de $\mathbb{R}$-álgebras $\varphi: A\to \mathbb{R}$, es decir, aplicaciones $\mathbb{R}$-lineales, multiplicativas, con $\varphi(1_A)=1$.

\begin{prop}
La aplicación
\begin{gather*}
    \mathrm{Alg}_{\mathbb{R}}(A,\mathbb{R}) \to \mathrm{Max}(A), \qquad \varphi \mapsto \ker(\varphi)
\end{gather*}
está bien definida: para cada $\varphi$, $\ker(\varphi)$ es un ideal maximal.
\end{prop}
\begin{proof}
Como $\varphi$ es sobreyectiva (salvo que sea la aplicación nula, que queda excluida por $\varphi(1)=1\neq 0$) sobre el cuerpo $\mathbb{R}$, el primer teorema de isomorfía da $A/\ker(\varphi) \cong \Im(\varphi) = \mathbb{R}$, que es un cuerpo. Por el Teorema~\ref{teo:correspondencia}, que $A/\ker(\varphi)$ sea cuerpo (sin ideales propios no nulos) equivale a que $\ker(\varphi)$ sea maximal.
\end{proof}

Esta aplicación en general \textbf{no} es ni inyectiva ni sobreyectiva; su estudio (para $A$ un anillo de polinomios o un cociente suyo) es precisamente el contenido del \textbf{Teorema de los ceros de Hilbert} (Nullstellensatz), que se enuncia y demuestra sobre un cuerpo algebraicamente cerrado (típicamente $\mathbb{C}$, no $\mathbb{R}$).

\begin{ejemplo}
Sea $A = \mathbb{R}[x]/\langle x^2 + 1 \rangle$. Como $x^2+1$ es irreducible en $\mathbb{R}[x]$ (no tiene raíces reales) y $\mathbb{R}[x]$ es un dominio de ideales principales, $\langle x^2+1\rangle$ es un ideal maximal, así que $A$ es un cuerpo; de hecho
\begin{gather*}
    A = \mathbb{R}[x]/\langle x^2+1\rangle \;\cong\; \mathbb{C}
\end{gather*}
(enviando $\bar x \mapsto i$). Entonces $\mathrm{Alg}_{\mathbb{R}}(A,\mathbb{R}) = \mathrm{Alg}_{\mathbb{R}}(\mathbb{C},\mathbb{R}) = \emptyset$: no existe ningún morfismo de $\mathbb{R}$-álgebras $\mathbb{C}\to \mathbb{R}$, porque tal morfismo debería enviar $i$ a un elemento $t\in \mathbb{R}$ con $t^2=\varphi(i)^2=\varphi(i^2)=\varphi(-1)=-1$, imposible en $\mathbb{R}$.
\end{ejemplo}

\emph{(Corrección: en el original se identificaba este cociente con $\mathbb{Q}(\sqrt{-1})$. El isomorfismo correcto es con $\mathbb{C}=\mathbb{R}(\sqrt{-1})$, ya que estamos cocientando $\mathbb{R}[x]$, no $\mathbb{Q}[x]$.)}

\subsection{El ejemplo de las funciones continuas}

Sea $X$ un espacio topológico y
\begin{gather*}
    C(X) = C(X,\mathbb{R}) = \{\text{funciones continuas } X\to \mathbb{R}\},
\end{gather*}
que es una $\mathbb{R}$-álgebra con las operaciones puntuales.

Para cada $x\in X$ podemos considerar
\begin{gather*}
    \mu_x = \{f \in C(X) : f(x)=0\},
\end{gather*}
que es el núcleo del morfismo de evaluación $\mathrm{ev}_x: C(X)\to \mathbb{R}$, $f\mapsto f(x)$, luego $\mu_x \in \mathrm{Max}(C(X))$. Esto define una aplicación bien definida
\begin{gather*}
    X \to \mathrm{Max}(C(X)), \qquad x \mapsto \mu_x.
\end{gather*}
En general esta aplicación no es biyectiva (por ejemplo si $X$ no es compacto puede haber "más" ideales maximales que puntos); el \textbf{Teorema de Gelfand-Kolmogorov} garantiza que sí lo es cuando $X$ es un espacio compacto y Hausdorff.

\begin{prop}
$X$ es conexo si y solo si $C(X)$ no tiene idempotentes no triviales, es decir, los únicos $e\in C(X)$ con $e^2=e$ son $e=0$ y $e=1$.
\end{prop}
\begin{proof}
Si $e^2=e$, entonces para cada $x$, $e(x)^2=e(x)$ en $\mathbb{R}$, luego $e(x)\in\{0,1\}$. Los conjuntos $e^{-1}(0)$ y $e^{-1}(1)$ son abiertos (preimágenes de abiertos por la aplicación continua $e$, ya que $\{0\}$ y $\{1\}$ son abiertos en el subespacio discreto $e(X)\subseteq\{0,1\}$) y forman una partición de $X$. Si $X$ es conexo, uno de los dos debe ser vacío, es decir $e\equiv 0$ o $e\equiv 1$. Recíprocamente, si $X=U\sqcup V$ con $U,V$ abiertos disjuntos no vacíos, la función $e=\mathbf 1_U$ (que vale $1$ en $U$ y $0$ en $V$) es continua e idempotente no trivial.
\end{proof}

\emph{(Corrección: el enunciado original decía, de forma poco precisa, "que $X$ sea conexo nos dice que $\varphi(x)$ es idempotente"; lo que realmente se afirma —y se demuestra arriba— es la equivalencia entre conexión de $X$ y ausencia de idempotentes no triviales en $C(X)$.)}

\subsection{Gérmenes de funciones}

Fijemos $x_0\in \mathbb{R}$. Aunque una función continua $f$ con $f(x_0)\neq 0$ no es invertible en $C(\mathbb{R})$ (puede anularse en otro punto), sí queremos poder "invertirla cerca de $x_0$". La idea es quedarse solo con el comportamiento de $f$ en un entorno arbitrariamente pequeño de $x_0$.

Formalmente, sea $\mathcal V_{x_0}$ el conjunto de entornos abiertos de $x_0$ en $\mathbb{R}$, y consideramos el conjunto de pares
\begin{gather*}
    \{(I, f) : I \in \mathcal V_{x_0},\ f\in C(I, \mathbb{R})\},
\end{gather*}
con la relación de equivalencia
\begin{gather*}
    (I,f)\, R_{x_0}\, (J,g) \ \sii\ \exists K\subseteq I\cap J \text{ entorno abierto de } x_0 \text{ tal que } f_{|K} = g_{|K}.
\end{gather*}
La clase de $(I,f)$ se llama el \textbf{germen} de $f$ en $x_0$, y se denota $[f]_{x_0}$ o $\overline{(I,f)}$. El conjunto de gérmenes
\begin{gather*}
    C_{x_0}(\mathbb{R}) := \Big(\bigsqcup_{I\in \mathcal V_{x_0}} C(I,\mathbb{R})\Big)\Big/ R_{x_0}
\end{gather*}
(el \textbf{límite directo}, o colímite, de los anillos $C(I,\mathbb{R})$ sobre los entornos de $x_0$) es de nuevo un anillo, de hecho una $\mathbb{R}$-álgebra, llamada el \textbf{anillo de gérmenes} (o \textbf{fibra}/\textbf{stalk}) de $C(-,\mathbb{R})$ en $x_0$.

Lo importante es que si $f(x_0)\neq 0$, entonces, por continuidad, $f$ no se anula en un entorno suficientemente pequeño de $x_0$, así que $1/f$ está bien definida y es continua en ese entorno, y por tanto $[f]_{x_0}$ es \textbf{invertible} en $C_{x_0}(\mathbb{R})$ — aunque $f$ misma no sea invertible en $C(\mathbb{R})$. Este proceso de "hacer invertible localmente" es un caso particular de un formalismo general llamado \textbf{localización}, que estudiamos a continuación: de hecho, $C_{x_0}(\mathbb{R})$ es (isomorfo a) la localización de $C(\mathbb{R})$ (o de $C(I)$ para cualquier $I\in\mathcal V_{x_0}$) en el conjunto multiplicativo $S_{x_0}=\{f : f(x_0)\neq 0\}$.

\section{Localización de anillos}

\begin{definicion}
Sea $A$ un anillo y $S\subseteq A$. Decimos que $S$ es un \textbf{conjunto multiplicativo} (o multiplicativamente cerrado) si:
\begin{enumerate}
    \item $1\in S$,
    \item $s,t\in S \Rightarrow st\in S$.
\end{enumerate}
(Convendremos además, salvo mención expresa de lo contrario, que $0\notin S$; en caso contrario la localización resulta ser el anillo nulo, como veremos.)
\end{definicion}

\begin{ejemplo}\
\begin{itemize}
    \item Si $A$ es un dominio y $A\neq 0$, $S=A\setminus\{0\}$ es multiplicativo (el producto de no nulos es no nulo, precisamente porque $A$ es dominio).
    \item Si $\mathfrak p \in \mathrm{Spec}(A)$, entonces $S=A\setminus \mathfrak p$ es multiplicativo: esta es justamente la propiedad (3) de la Proposición~\ref{prop:primo-equivalencias}.
    \item Si $f\in A$, $S=\{f^n : n\geq 0\}$ (con $f^0=1$) es multiplicativo, siempre que $f$ no sea nilpotente (si algún $f^n=0$, entonces $0\in S$).
\end{itemize}
\end{ejemplo}

\subsection{Construcción de $AS^{-1}$}

Sea $A$ un anillo y $S\subseteq A$ multiplicativo. Consideramos el producto cartesiano $A\times S$ (pensando en $(a,s)$ como el futuro "$a/s$") y la relación
\begin{gather*}
    (a,s)\, R_S\, (b,t) \ \overset{\mathrm{def}}{\sii}\ \exists u\in S \ : \ u(at-sb) = 0.
\end{gather*}

\begin{prop}
$R_S$ es una relación de equivalencia.
\end{prop}
\begin{proof}
\emph{Reflexiva}: $(a,s)R_S(a,s)$ pues $1\cdot(as-sa)=0$, con $1\in S$.

\emph{Simétrica}: si $u(at-sb)=0$ entonces $u(bs-ta)=-u(at-sb)=0$, luego $(b,t)R_S(a,s)$.

\emph{Transitiva}: supongamos $(a,s)R_S(b,t)$ y $(b,t)R_S(c,r)$, es decir, existen $u,v\in S$ con
\begin{gather*}
    u(at-sb)=0, \qquad v(br-tc)=0.
\end{gather*}
Multiplicando la primera por $vr$ y la segunda por $us$ y sumando:
\begin{gather*}
    uvr(at-sb) + usv(br-tc) = uv\,t\,(ar-sc) = 0
\end{gather*}
(los términos con $b$ se cancelan: $-uvrsb+uvsbr=0$). Como $uvt\in S$ (producto de elementos de $S$), esto prueba $(a,s)R_S(c,r)$.
\end{proof}

Denotamos $AS^{-1} := (A\times S)/R_S$, y la clase de $(a,s)$ se escribe $\dfrac{a}{s}$.

Definimos las operaciones
\begin{gather*}
    \frac{a}{s} + \frac{b}{t} := \frac{at+sb}{st}, \qquad\qquad \frac{a}{s} \cdot \frac{b}{t} := \frac{ab}{st},
\end{gather*}
con $0_{AS^{-1}} = \dfrac{0}{1}$ y $1_{AS^{-1}} = \dfrac{1}{1}$.

\emph{(Corrección: en el original la suma se escribía $\frac{a}{s}+\frac{b}{r}=\frac{as+sb}{sr}$, con un error tipográfico —debía ser $at+sb$ sobre el denominador común $st$, no $as+sb$—.)}

Ambas operaciones están bien definidas (no dependen de los representantes elegidos): esto se comprueba usando exactamente la definición de $R_S$, de forma completamente análoga al cálculo que hicimos para el producto en un anillo cociente.

\begin{prop}
$(AS^{-1}, +, \cdot)$ es un anillo conmutativo.
\end{prop}

Tenemos el \textbf{morfismo canónico de localización}
\begin{gather*}
    \tau: A \to AS^{-1}, \qquad a \mapsto \frac{a}{1},
\end{gather*}
que es morfismo de anillos.

\begin{prop}
Si $A$ es un dominio (y $0\notin S$), entonces $\tau$ es inyectivo.
\end{prop}
\begin{proof}
\begin{gather*}
    \frac{a}{1} = \frac{0}{1} \ \sii\ \exists s \in S \ : \ s(a\cdot 1 - 1\cdot 0) = 0 \ \sii\ \exists s\in S,\ sa=0.
\end{gather*}
Como $A$ es dominio y $s\neq 0$ (pues $0\notin S$), $sa=0 \Rightarrow a=0$. Luego $\ker(\tau)=0$.
\end{proof}

En general (si $A$ no es dominio, o si $S$ contiene divisores de cero) $\tau$ puede no ser inyectivo; de hecho, si $0\in S$, entonces para cualquier $a,b$ y $s=0\in S$: $s(a\cdot 1-1\cdot b)=0$ automáticamente, así que todos los pares son equivalentes y $AS^{-1}$ es el anillo nulo.

\begin{teo}[Propiedad universal de la localización]
Sea $\varphi: A \to B$ un morfismo de anillos tal que $\varphi(S) \subseteq U(B)$ (las imágenes de los elementos de $S$ son invertibles en $B$). Entonces existe un único morfismo $\tilde{\varphi}: AS^{-1}\to B$ tal que $\tilde\varphi\circ \tau = \varphi$, es decir, que hace conmutar
\begin{gather*}
    \xymatrix{
        A \ar[d]_{\tau} \ar[r]^{\varphi} & B\\
        AS^{-1} \ar@{-->}[ru]_{\tilde{\varphi}}
    }
\end{gather*}
Explícitamente, $\tilde\varphi(a/s) = \varphi(a)\varphi(s)^{-1}$.
\end{teo}

\subsection{Funtorialidad}

Si tenemos un morfismo $\varphi:A\to B$ (sin hipótesis adicionales) y fijamos $S\subseteq A$ multiplicativo, entonces $T:=\varphi(S)\subseteq B$ es multiplicativo. Componiendo $\varphi$ con $\tau_B: B\to BT^{-1}$ obtenemos un morfismo $A\to BT^{-1}$ que envía $S$ a unidades (las imágenes de $S$ están en $T$, que se vuelve invertible en $BT^{-1}$); por la propiedad universal, se factoriza de forma única a través de $AS^{-1}$, dando el \textbf{morfismo localizante}
\begin{gather*}
    \xymatrix{
        A \ar[d]_{\tau_A} \ar[r]^{\varphi} & B \ar[d]^{\tau_B}\\
        AS^{-1} \ar[r]_{\varphi_S} & BT^{-1}
    }
\end{gather*}

Además, cualquier morfismo de anillos $\varphi: A\to B$ (no solo entre localizaciones) induce una aplicación en sentido contrario entre espectros,
\begin{gather*}
    \varphi^{\#}: \mathrm{Spec}(B) \to \mathrm{Spec}(A), \qquad \mathfrak q \mapsto \varphi^{-1}(\mathfrak q),
\end{gather*}
que está bien definida porque la preimagen de un ideal primo por un morfismo de anillos es un ideal primo (ejercicio directo usando la caracterización (2) de la Proposición~\ref{prop:primo-equivalencias}: $A/\varphi^{-1}(\mathfrak q)$ se inyecta en $B/\mathfrak q$, que es dominio, luego $A/\varphi^{-1}(\mathfrak q)$ también lo es).

\subsection{Extensión y contracción de ideales}

Sea $\tau: A\to AS^{-1}$ la localización.
\begin{itemize}
    \item Si $I\leq A$, su \textbf{extensión} es el ideal de $AS^{-1}$ generado por $\tau(I)$:
    \begin{gather*}
        I^e := \langle \tau(I)\rangle_{AS^{-1}} = \Big\{\frac{a}{s} \ :\ a\in I,\ s\in S\Big\}.
    \end{gather*}
    \item Si $J\leq AS^{-1}$, su \textbf{contracción} es
    \begin{gather*}
        J^c := \tau^{-1}(J) \leq A.
    \end{gather*}
\end{itemize}

\emph{(Corrección: en el original se usaba el mismo símbolo "$e$" tanto para la extensión como para la contracción, escribiendo "$J^e=\tau^{-1}(J)$"; la notación estándar, que usamos aquí, distingue extensión ($I^e$) de contracción ($J^c$).)}

\begin{prop}
Con las notaciones anteriores:
\begin{enumerate}
    \item $I^e = \{a/s : a\in I,\ s\in S\}$ (ya visto).
    \item Si $I\cap S\neq \emptyset$, entonces $I^e = AS^{-1}$ (el ideal total).
    \item $(J^c)^e = J$ para todo ideal $J\leq AS^{-1}$, y $I \subseteq (I^e)^c$ para todo ideal $I\leq A$ (en general sin igualdad).
\end{enumerate}
\end{prop}
\begin{proof}[Idea]
(2): si $u\in I\cap S$, entonces $1=u/u \in I^e$ (pues $u\in I$, $u\in S$), luego $I^e$ contiene una unidad y es el ideal total.

(3): $(J^c)^e\subseteq J$ es claro; para la otra inclusión, dado $a/s\in J$, se tiene $\tau(a)=s\cdot(a/s)\in J$ (usando que $s/1$ es invertible en $AS^{-1}$, con inverso $1/s$), luego $a\in J^c=\tau^{-1}(J)$ y $a/s = a/s \in (J^c)^e$.

Para $I\subseteq (I^e)^c$: si $a\in I$, entonces $\tau(a)=a/1 \in I^e$ por definición, luego $a\in \tau^{-1}(I^e) = (I^e)^c$.

La inclusión $I\subseteq (I^e)^c$ puede ser estricta: basta que $I$ contenga elementos que, tras localizar, "se conviertan" en elementos de un ideal más pequeño (ver el ejemplo siguiente).
\end{proof}

\begin{ejemplo}
Sea $A=\mathbb{Z}$ y $S=\{5^n : n\geq 0\}$ (así $AS^{-1}=\mathbb{Z}[1/5]$). En $AS^{-1}$ tenemos la igualdad
\begin{gather*}
    \frac{2}{5} = \frac{10}{25}
\end{gather*}
(ambas representan la misma clase, ya que $1\cdot(2\cdot 25 - 5\cdot 10)=0$).

Consideremos el ideal $I = 10\mathbb{Z} \leq \mathbb{Z}$. Se tiene $10\in I$, pero $2\notin I$. Sin embargo, usando la representación $\frac{2}{5}=\frac{10}{25}$ con $10\in I$, obtenemos que $\frac{2}{5}\in I^e$ (por la caracterización (1) anterior, tomando $a=10\in I$, $s=25\in S$), \textbf{a pesar de que su "numerador reducido" $2$ no esté en $I$}.

Esto ilustra un punto delicado: para decidir si una fracción $a/s$ pertenece a $I^e$ \textbf{no} basta mirar si el numerador de la representación "más simple" está en $I$; hay que comprobar si \emph{existe alguna} representación equivalente con numerador en $I$. En este ejemplo concreto, de hecho $I^e = \langle \tau(10)\rangle = \langle 2\rangle_{AS^{-1}}$ (porque $5$ es invertible en $AS^{-1}$, luego $10=2\cdot 5$ y $2$ son asociados), mientras que como ideales de $\mathbb{Z}$, $\langle 10\rangle \neq \langle 2\rangle$: la extensión "olvida" los factores de $I$ que caen en $S$.
\end{ejemplo}

\emph{(Corrección: el ejemplo original terminaba abruptamente con la igualdad mal escrita "$\frac{-2}{5}=\frac{-10}{2}$" —que es falsa: $-10/2=-5\neq -2/5$— seguida de la observación sobre el ideal $10\mathbb{Z}$ sin conectar ambas cosas. Se ha reconstruido el ejemplo con la igualdad correcta $\frac{2}{5}=\frac{10}{25}$ y se ha completado la conclusión que se pretendía ilustrar: la extensión de un ideal por localización no se calcula mirando únicamente los numeradores "reducidos".)}